\scanpage{42}
\mspage{41}

En particulier, si $G_0,G'_0$ satisfont aux conditions de la proposition, alors $\operatorname{Hom}(G_0,G'_0)\ne0$, i.e. $G_0$ et $G'_0$ sont isogènes, ssi $\pi=\pi'$. \deleted{\unsure{Si}} S'il en est ainsi, alors \deleted{\unsure{on a}} on a
\[
\left\{\begin{array}{ll}
\text{si }\nu\unsure{\leq|<}\nu'&\deleted{G_0=G'_0},\quad\operatorname{Hom}_A(G_0,G'_0)\overset\sim\longrightarrow\operatorname{Hom}(\Omega',\Omega)\\
\text{si }\nu>\nu'&\operatorname{Hom}_A(G_0,G'_0)\simeq\pi\operatorname{Hom}(\Omega',\Omega)
\end{array}\right.
\]

\uline{Re NB} Revenons au cas général où on ne suppose pas néc. $\dim G_0=1$, \uuline{ni} \unsure{sa} hauteur $=h$. Classifions à isogénie près \deleted{\unsure{les}} ces $G_0$, donc les $M_\nu$ $(0\leq\nu\leq f-1)$ deviennent des vectoriels \deleted{de dim.}\inserted{sur le} corps $K$, isomorphes par les $\varphi_\nu$ $(0\leq\nu\leq f-2)$. Alors
\deleted{\unsure{on est en présence de}}
\[
M\simeq\Omega\otimes_{\mathbb Q_p}W=\Omega\otimes_K(K\otimes_{\mathbb Q_p}K_0)
\]
\margin{$K_0$ corps des fractions de $W$}
et $F$ est déterminé par la connaissance de $\varphi_{f-1}$, qui est un endom de $M$, soit $\varpi$, de sorte qu'on a
\[
F_M^f=\varpi\otimes_{\mathbb Q_p}\operatorname{id}_{K_0}=\varpi\otimes_K(K\otimes_{\mathbb Q_p}K_0).
\]
La condition ``d'effectivité'' en termes de $\varpi$ est
\[
\left\{\begin{array}{l}
\varpi\ \text{laisse invariant un réseau}\\
p^f\varpi^{-1}\ \text{laisse invariant un réseau}
\end{array}\right.
\]
